\chapter{Podsumowanie i~wnioski}
\label{chap:wnioski}

W~podsumowaniu w~zwięzły sposób przypomina się cel pracy i~ocenia, czy oraz
w~jakim stopniu został on osiągnięty. Wnioski powinny wynikać bezpośrednio
z~przeprowadzonych obliczeń, analiz lub badań -- nie wprowadza się tu
nowych informacji ani odwołań do literatury, których wcześniej nie
omówiono.

Wnioski wygodnie jest sformułować w~punktach, np.:
\begin{enumerate}
  \item Projektowe obciążenie cieplne budynku wynosi \qty{8.4}{kW}, co
        pozwala na zastosowanie pompy ciepła o~mocy nominalnej \dots
  \item Roczny koszt eksploatacji wariantu z~pompą ciepła jest o~\dots\%
        niższy niż wariantu z~kotłem gazowym.
  \item Największy wpływ na wynik analizy ekonomicznej ma \dots
\end{enumerate}

Jeżeli praca ma charakter projektowy, należy podać zalecenia wynikające
z~obliczeń. Na końcu można wskazać kierunki dalszych prac, np.\
rozszerzenie analizy o~symulację całoroczną albo pomiary na obiekcie
rzeczywistym.
